\subsection{Relative Assessment Prompt}

\texttt{---Role---\\
\\
You are a helpful assistant responsible for grading two answers to a question that are provided by two different people.\\
\\
---Goal---\\
\\
Given a question and two answers (Answer 1 and Answer 2), assess which answer is better according to the following measure:\\
\\
\{criteria\}\\
\\
Your assessment should include two parts:\\
- Winner: either 1 (if Answer 1 is better) and 2 (if Answer 2 is better) or 0 if they are fundamentally similar and the differences are immaterial.\\
- Reasoning: a short explanation of why you chose the winner with respect to the measure described above.\\
\\
Format your response as a JSON object with the following structure:\\
\{\{\\
    "winner": <1, 2, or 0>,\\
    "reasoning": "Answer 1 is better because <your reasoning>."\\
\}\}\\
\\
---Question---\\
\\
\{question\}\\
\\
---Answer 1---\\
\\
\{answer1\}\\
\\
---Answer 2---\\
\\
\{answer2\}\\
\\
Assess which answer is better according to the following measure:\\
\\
\{criteria\}\\
\\
Output:}

\subsection{Relative Assessment Metrics}

\texttt{CRITERIA = \{\\
    "comprehensiveness": "How much detail does the answer provide to cover all the aspects and details of the question? A comprehensive answer should be thorough and complete, without being redundant or irrelevant. For example, if the question is 'What are the benefits and drawbacks of nuclear energy?', a comprehensive answer would provide both the positive and negative aspects of nuclear energy, such as its efficiency, environmental impact, safety, cost, etc. A comprehensive answer should not leave out any important points or provide irrelevant information. For example, an incomplete answer would only provide the benefits of nuclear energy without describing the drawbacks, or a redundant answer would repeat the same information multiple times.",\\
    "diversity": "How varied and rich is the answer in providing different perspectives and insights on the question? A diverse answer should be multi-faceted and multi-dimensional, offering different viewpoints and angles on the question. For example, if the question is 'What are the causes and effects of climate change?', a diverse answer would provide different causes and effects of climate change, such as greenhouse gas emissions, deforestation, natural disasters, biodiversity loss, etc. A diverse answer should also provide different sources and evidence to support the answer. For example, a single-source answer would only cite one source or evidence, or a biased answer would only provide one perspective or opinion.",\\
    "directness": "How specifically and clearly does the answer address the question? A direct answer should provide a clear and concise answer to the question. For example, if the question is 'What is the capital of France?', a direct answer would be 'Paris'. A direct answer should not provide any irrelevant or unnecessary information that does not answer the question. For example, an indirect answer would be 'The capital of France is located on the river Seine'.",\\
    "empowerment": "How well does the answer help the reader understand and make informed judgements about the topic without being misled or making fallacious assumptions. Evaluate each answer on the quality of answer as it relates to clearly explaining and providing reasoning and sources behind the claims in the answer."\\
\}}